\documentclass[a4paper]{article}

\usepackage[utf8]{inputenc}
\usepackage[T2A]{fontenc}
\usepackage[russian,english]{babel}
\usepackage{amsmath}
\usepackage{graphicx}
\usepackage{caption}
\usepackage{subcaption}
\usepackage{float}
\usepackage[ruled,vlined,linesnumbered]{algorithm2e} % стиль "красивого блока"
\usepackage{amsmath,amssymb}
\usepackage{xcolor}
\usepackage[colorinlistoftodos]{todonotes}
\usepackage{booktabs}
\usepackage{siunitx}
\sisetup{
  detect-weight=true,
  detect-inline-weight=math,
  table-number-alignment=center,
  separate-uncertainty=true
}
\setlength{\tabcolsep}{4pt}

% (опционально) чуть приятнее шрифт и отступы внутри алгоритмов
\SetAlgoInsideSkip{smallskip}
\SetAlgoSkip{smallskip}
\DontPrintSemicolon % если не хочешь ; в конце строк (в примере как на скрине)

% (опционально) серые "комменты справа" как на скрине
\SetCommentSty{itshape\color{gray}}
\SetKwComment{Comment}{$\triangleright$ }{}  % значок ▶ как на примере
\SetKwInput{KwIn}{Input}
\SetKwInput{KwOut}{Output}



\title{Результаты экспериментов по планированию с GCS}
\date{}
\begin{document}
\maketitle


\section{Разбиение на выпуклые регионы}
В предыдущих экспериментах начальные точки для построения IRIS-регионов генерировались равномерно и случайно внутри ограничивающего объёма сцены (bounding box). Такой подход приводил к недостаточному покрытию областей со сложной геометрией, в частности зон вблизи полок стеллажей, из-за чего в этих областях формировалось меньшее число выпуклых регионов.

Поскольку при планировании движений манипулятора важны траектории, проходящие в непосредственной близости к полкам, нужно более плотное покрытие таких зон выпуклыми оболочками.

Для этого при генерации начальных точек для построения выпуклых оболочек был введет параметр $shelf\_ratio$, который отвечает за то, какой процент начальных точек генерируется в области сложной геометрии (в нашем случае полок на стеллажах). Во всех экспериментах ниже использовался $shelf\_ratio = 0.3$


\section{Описание сценариев}
Оценка конечных траекторий движения производилась по двум параметрам: длина пути и время выполнения. Для проведения экспериментов было выбрана два сценария:
\begin{enumerate}
    \item Планирование пути между максимально удаленными точками.
    \item Планирование пути в сложных геометриях (в нашем случае рядом с полками).
\end{enumerate}
Как и раньше, использовались 3 сцены (kuka iiwa с одной полкой, двумя и тремя). Сравнение производилось по 10, 20, 30 и 50 выпуклым регионам. То есть
для каждой сцены строилось 10, 20, 30 и 50 выпуклых регионов.

\section{Алгоритм построения экспериментов}

При попытке напрямую применить алгоритм планирования на основе GCS (реализация Drake) выяснилось, что при количестве выпуклых регионов более 30 задача оптимизации не сходится из-за слишком большого размера задачи.
Поэтому был использован двухэтапный подход. Сначала с помощью алгоритма поиска в ширину (BFS) находился кратчайший путь в графе выпуклых регионов между регионом, содержащим начальную конфигурацию,
и регионом, содержащим целевую конфигурацию. Затем алгоритм планирования GCS запускался только на подграфе регионов, соответствующем найденному пути, что существенно снижало вычислительную сложность и позволяло получить решение.



\begin{algorithm}[H]
\caption{Методика проведения экспериментов с GCS планированием}
\label{alg:gcs_experiments}

\KwIn{Типы сцен $\mathcal{S}$, количества регионов $N = \{10, 20, 30, 50\}$, число попыток $T$, параметр $shelf\_ratio = 0.3$}
\KwOut{Статистика по длинам путей и времени решения}

\BlankLine

\ForEach{$scene \in \mathcal{S}$}{
    \ForEach{$n \in N$}{
        \ForEach{$trial = 1$ \KwTo $T$}{
            $\mathcal{R} \gets$ \texttt{BuildIRISRegions}$(scene, n, shelf\_ratio)$ \;
            \texttt{SaveRegions}$(\mathcal{R}, \text{``regions/''})$ \;
            
            \BlankLine
            \tcp{Experiment 1: Distant points}
            $\mathcal{Q} \gets$ \texttt{SampleCollisionFree}$(scene, 200)$ \;
            $\mathcal{Q}_{covered} \gets \{q \in \mathcal{Q} : \exists r \in \mathcal{R}, q \in r\}$ \;
            $(q_{start}, q_{goal}) \gets$ \texttt{FindDistantPair}$(\mathcal{Q}_{covered})$ \;
            \If{$\|q_{start} - q_{goal}\| \geq d_{min}$}{
                $\mathcal{R}_{path} \gets$ \texttt{BFS}$(q_{start}, q_{goal}, \mathcal{R})$ \;
                $path \gets$ \texttt{GCS-Solve}$(q_{start}, q_{goal}, \mathcal{R}_{path})$ \;
                \texttt{RecordResult}$(path, \text{``distant''})$ \;
            }
            
            \BlankLine
            \tcp{Experiment 2: Shelf-to-shelf paths}
            $\mathcal{P} \gets$ \texttt{GetShelfPositions}$(scene)$ \;
            \ForEach{$(p_i, p_j) \in \mathcal{P} \times \mathcal{P}$, $i < j$}{
                $q_i \gets$ \texttt{SolveIK}$(p_i)$, $q_j \gets$ \texttt{SolveIK}$(p_j)$ \;
                \If{$q_i \in \bigcup \mathcal{R}$ \textbf{and} $q_j \in \bigcup \mathcal{R}$}{
                    $\mathcal{R}_{path} \gets$ \texttt{BFS}$(q_i, q_j, \mathcal{R})$ \;
                    $path \gets$ \texttt{GCS-Solve}$(q_i, q_j, \mathcal{R}_{path})$ \;
                    \texttt{RecordResult}$(path, \text{``shelf''})$ \;
                }
            }
        }
    }
}

\BlankLine
\Return Агрегированная статистика по всем экспериментам
\end{algorithm}

\textbf{Ключевые особенности алгоритма:}
\begin{itemize}
    \item \textbf{Смешанная стратегия генерации:} 30\% начальных точек для IRIS размещаются вблизи полок ($shelf\_ratio = 0.3$), остальные 70\% -- случайно в пространстве конфигураций.
    
    \item \textbf{Двухэтапное планирование:} Для сцен с $n > 30$ регионами сначала применяется BFS для поиска последовательности регионов, затем GCS решает задачу оптимизации только на этом подграфе.
    
    \item \textbf{Эксперимент ``distant'':} Оценивает качество планирования между максимально удалёнными точками в $C$-пространстве ($d_{min} = 2.0$).
    
    \item \textbf{Эксперимент ``shelf'':} Проверяет способность планировщика находить траектории в сложной геометрии (между позициями на полках стеллажей).
    
    \item \textbf{Проверка покрытия:} Обе конечные точки должны принадлежать хотя бы одному выпуклому региону, иначе эксперимент пропускается.
\end{itemize}


\section{Результаты экспериментов и анализ}

Результаты экспериментов по планированию траекторий с использованием GCS представлены в таблицах~\ref{tab:distant_results} и~\ref{tab:shelf_results}. Эксперименты проводились для трёх сцен различной сложности с разным количеством выпуклых регионов (10, 20, 30, 50).

\subsection{Эксперимент 1: Планирование между удалёнными точками}

\begin{table}[H]
\centering
\caption{Результаты планирования между максимально удалёнными точками в C-пространстве.}
\label{tab:distant_results}
\begin{tabular}{l
                S[table-format=2.0]
                S[table-format=1.0]/S[table-format=1.0]
                S[table-format=2.2]
                S[table-format=1.4]}
\toprule
\textbf{Scene} & \textbf{Seeds} & \textbf{Success} & {\textbf{Path Length}} & {\textbf{Solve Time (s)}} \\
\midrule
Single Shelf   & 10 & 3/3 & 10.31 \pm 0.20 & 0.007 \\
Single Shelf   & 20 & 3/3 & 11.29 \pm 1.00 & 0.012 \\
Single Shelf   & 30 & 3/3 & 10.73 \pm 0.62 & 0.008 \\
Single Shelf   & 50 & 3/3 & 11.03 \pm 0.46 & 0.007 \\
\addlinespace
Three Shelves  & 10 & 3/3 & 9.77 \pm 0.51  & 0.019 \\
Three Shelves  & 20 & 3/3 & 10.52 \pm 1.59 & 0.021 \\
Three Shelves  & 30 & 3/3 & 10.30 \pm 0.42 & 0.013 \\
Three Shelves  & 50 & 3/3 & 10.99 \pm 0.76 & 0.027 \\
\addlinespace
Two Shelves    & 10 & 3/3 & 12.93 \pm 1.25 & 0.029 \\
Two Shelves    & 20 & 3/3 & 11.81 \pm 1.44 & 0.010 \\
Two Shelves    & 30 & 3/3 & 11.00 \pm 0.65 & 0.014 \\
Two Shelves    & 50 & 3/3 & 11.06 \pm 1.22 & 0.008 \\
\bottomrule
\end{tabular}
\end{table}

Эксперименты по планированию между максимально удалёнными точками показали \textbf{100\% успешность} для всех конфигураций. Длины путей находятся в диапазоне 9.77--12.93, что говорит о стабильности алгоритма. Время решения очень мало (0.007--0.029 с), что подтверждает эффективность двухэтапного подхода с использованием BFS для предварительной фильтрации регионов.

\subsection{Эксперимент 2: Планирование между позициями на полках}

\begin{table}[H]
\centering
\caption{Результаты планирования траекторий между позициями на полках стеллажей.}
\label{tab:shelf_results}
\begin{tabular}{l
                S[table-format=2.0]
                S[table-format=2.0]/S[table-format=2.0]
                S[table-format=1.0]
                S[table-format=1.2]
                S[table-format=1.4]}
\toprule
\textbf{Scene} & \textbf{Seeds} & \textbf{Success} & {\textbf{Rate (\%)}} & {\textbf{Path Length}} & {\textbf{Solve Time (s)}} \\
\midrule
Single Shelf   & 10 & 3/3   & 100 & 0.48 \pm 0.06 & 0.010 \\
Single Shelf   & 20 & 1/3   & 33  & 0.07 \pm 0.13 & 0.002 \\
Single Shelf   & 30 & 1/3   & 33  & 0.18 \pm 0.31 & 0.002 \\
Single Shelf   & 50 & 1/3   & 33  & 0.07 \pm 0.13 & 0.001 \\
\addlinespace
Three Shelves  & 10 & 3/21  & 14  & 0.11 \pm 0.31 & 0.002 \\
Three Shelves  & 20 & 11/21 & 52  & 0.56 \pm 0.82 & 0.001 \\
Three Shelves  & 30 & 12/21 & 57  & 0.70 \pm 1.08 & 0.003 \\
Three Shelves  & 50 & 15/21 & 71  & 1.07 \pm 1.19 & 0.010 \\
\addlinespace
Two Shelves    & 10 & 0/15  & 0   & {---}         & {---} \\
Two Shelves    & 20 & 1/15  & 7   & 0.01 \pm 0.05 & 0.000 \\
Two Shelves    & 30 & 10/15 & 67  & 1.02 \pm 1.10 & 0.006 \\
Two Shelves    & 50 & 0/15  & 0   & {---}         & {---} \\
\bottomrule
\end{tabular}
\end{table}

Результаты экспериментов по планированию между позициями на полках значительно более неоднородны и выявляют несколько важных закономерностей:

\begin{enumerate}
    \item \textbf{Влияние количества регионов на успешность:} Для сцены с тремя полками наблюдается монотонное увеличение успешности с 14\% (10 регионов) до 71\% (50 регионов). Это подтверждает, что более плотное покрытие пространства конфигураций выпуклыми регионами критично для сложных сценариев.
    
    \item \textbf{Проблемы с покрытием в сложных геометриях:} Для сцены Single Shelf успешность падает с 100\% (10 регионов) до 33\% (20--50 регионов), что может указывать на деградацию покрытия критических областей при увеличении общего числа регионов или на проблемы с решением обратной кинематики.
    
    \item \textbf{Нестабильность для сцены Two Shelves:} Результаты показывают крайне низкую успешность (0--67\%), причём для 50 регионов успешность равна нулю. Это говорит о недостаточном покрытии узких проходов между полками или о геометрических особенностях сцены, затрудняющих решение IK.
    
    \item \textbf{Длины путей:} Успешные траектории имеют существенно меньшую длину (0.01--1.07) по сравнению с экспериментом на удалённых точках, что логично, так как позиции на полках находятся относительно близко друг к другу.
\end{enumerate}

\section{Дальнейшее направление работы}

\textbf{Идея 1}:
рассмотреть планирование траекторий в GCS с учетом геометрических свойств выпуклых регионов.
Помимо стандартной минимизации длины пути вводится дополнительный штраф за прохождение через узкие области конфигурационного пространства.
так мы сможем предпочитать более “безопасные” траектории в сценах с узкими проходами.  \\
\\
\textbf{Идея 2}:
использовать простые графовые характеристики (степень вершины, точки сочленения, мосты) для модификации задачи планирования в GCS. Это позволяет влиять на выбор пути еще до запуска оптимизации.\\
\\
\textbf{Идея 3}:
планирование траекторий манипулятора, разделенное на фазы движения, например подход к полке и перемещение в свободном пространстве. Для разных фаз вводятся различные ограничения и штрафы в задаче GCS, что позволяет учитывать специфику движения вблизи препятствий.

\end{document}
